\documentclass[prd,twocolumn,floatfix,nofootinbib]{revtex4}
\usepackage{amssymb}
\usepackage{amsmath}
\usepackage{graphicx,color,dcolumn,booktabs,bm}
\usepackage{longtable,lscape}
\usepackage{txfonts}
\usepackage{overpic}
\usepackage{indentfirst}
\usepackage{cases}
\usepackage{multirow}
\usepackage{epstopdf}
\usepackage{ulem}
\usepackage[colorlinks,
            citecolor=blue,
            anchorcolor=red,
            menucolor=red,
            linkcolor=red,
            filecolor=red,
            runcolor=red,
            urlcolor=blue,
            frenchlinks=false]{hyperref}
%\newtheorem{theorem}{Theorem}
%\newtheorem{remark}[theorem]{Remark}
%\newtheorem{solution}[theorem]{Solution}
%\newtheorem{summary}[theorem]{Summary}
%\newenvironment{proof}[1][Proof]{\noindent\textbf{#1.} }{\ \rule{0.5em}{0.5em}}
%\usepackage{tikz}
\setcounter{MaxMatrixCols}{10}
\def\lrpartial{\buildrel\leftrightarrow\over\partial}
\graphicspath{{Figures/}} %
\allowdisplaybreaks

\linespread{1.1}
\begin{document}

\renewcommand{\tabcolsep}{0.72cm}
\renewcommand{\arraystretch}{1.3}
\begin{table*}[!htbp]
    \caption{Predicted spectra of ground-state and mixing heavy baryons. The experimental data $M_{exp}$ takes average of isospin.
    Magnetic moments and charge radii are organized in the order of $I_{3}=\frac{1}{2},\,-\frac{1}{2}$ for $I=\frac{1}{2}$.}
    \label{tab:mixingbaryon}
    \begin{tabular}{ccccccc}
        \hline\hline
        {\rm State}         &{\rm Eigenvector}  &$R_{0}$    &$M_{bag}$  &$M_{exp}$  &$\mu_{bag}/\mu_{p}$    &$r_{E}$ \\ \hline
        \multirow{2}{*}{$\left(\Xi_{c}^{\prime},\ \Xi_{c}\right)$}
                            &(0.05, 1.00)       &4.97	    &2.437	    &2.470      &0.13, 0.18             &0.64, 0.33 \\ 
                            &(-1.00, 0.05)      &4.97	    &2.546	    &2.579      &0.24, -0.43            &0.64, 0.33 \\ 
        $\Xi_{c}^{\ast}$ 	&1.00               &5.06	    &2.636	    &2.646      &0.58, -0.39            &0.65, 0.33 \\ 
        \multirow{2}{*}{$\left(\Xi_{b}^{\prime},\ \Xi_{b}\right)$}
                            &(0.01, 1.00)       &4.73	    &5.807	    &5.795      &-0.04, -0.03           &0.30, 0.61 \\ 
                            &(-1.00, 0.01)      &4.73	    &5.956	    &5.935      &0.27, -0.35            &0.30, 0.61 \\ 
        $\Xi_{b}^{\ast}$ 	&1.00               &4.79	    &5.991	    &5.954      &0.34, -0.58            &0.31, 0.62 \\ 
        \multirow{2}{*}{$\left(\Xi_{bc}^{\prime},\ \Xi_{bc}\right)$}
                            &(-0.39, -0.92)     &4.23	    &7.015	    &-          &0.53, -0.12            &0.58, 0.19 \\ 
                            &(-0.92, 0.39)      &4.23	    &6.955	    &-          &-0.08, 0.03            &0.58, 0.19 \\ 
        $\Xi_{bc}^{\ast}$ 	&1.000              &4.31	    &7.044	    &-          &0.69, -0.13            &0.59, 0.20 \\ 
        \multirow{2}{*}{$\left(\Omega_{bc}^{\prime},\ \Omega_{bc}\right)$}
                            &(-0.40, -0.92)     &4.31	    &7.117	    &-          &-0.07                  &0.15 \\ 
                            &(-0.92, 0.40)      &4.31	    &7.068	    &-          &0.02                   &0.15 \\ 
        $\Omega_{bc}^{\ast}$&1.000              &4.37	    &7.143	    &-          &-0.08                  &0.15 \\ 
        \hline\hline
    \end{tabular}
\end{table*}

\renewcommand{\tabcolsep}{0.62cm}
\renewcommand{\arraystretch}{1.3}
\begin{table*}[!htbp]
    \caption{Predicted spectra of singly heavy tetraquarks $nn\bar{s}\bar{c}$ and $nn\bar{s}\bar{b}$. 
    Magnetic moments and charge radii are organized in the order of $I_{3}=1,\,0,\,-1$ for $I=1$.}
    \label{tab:1heavytetraquark}
    \begin{tabular}{ccccccc}
        \hline\hline
        {\rm State}     &$J^{P}$    &{\rm Eigenvector}      &$R_{0}$    &$M_{bag}$  &$\mu_{bag}/\mu_{p}$    &$r_{E}$ \\ \hline
        \multirow{2}{*}{${(nn\bar{s}\bar{c})}^{I=1}$}
                        &\multirow{2}{*}{$0^{+}$}
                                    &(0.55, 0.84)           &5.56	    &3.219	    &-                      &0.88, 0.37, 0.71 \\ 
                        &           &(-0.84, 0.55)          &5.56	    &2.781	    &-                      &0.88, 0.37, 0.71 \\ 
                        &\multirow{3}{*}{$1^{+}$}
                                    &(0.80, 0.59, 0.10)     &5.56	    &3.001	    &1.27, 0.56, -0.15      &0.88, 0.37, 0.71 \\ 
                        &           &(0.25, -0.49, 0.84)    &5.56	    &3.154	    &0.37, 0.08, -0.20      &0.88, 0.37, 0.71 \\ 
                        &           &(0.54, -0.65, -0.53)   &5.56	    &2.848	    &0.62, 0.02, -0.59      &0.88, 0.37, 0.71 \\ 
                        &$2^{+}$    &1.00                   &5.64	    &3.075	    &1.53, 0.45, -0.63      &0.89, 0.37, 0.72 \\ 
        \multirow{2}{*}{${(nn\bar{s}\bar{c})}^{I=0}$}
                        &\multirow{2}{*}{$0^{+}$}
                                    &(0.63, 0.77)           &5.56	    &2.934	    &-                      &0.37 \\ 
                        &           &(-0.78, 0.63)          &5.56	    &2.526	    &-                      &0.37 \\ 
                        &\multirow{3}{*}{$1^{+}$}
                                    &(0.77, 0.06, 0.64)     &5.56	    &2.897	    &0.19                   &0.37 \\ 
                        &           &(-0.23, -0.90, 0.36)   &5.56	    &3.056	    &0.44                   &0.37 \\ 
                        &           &(0.60, -0.43, -0.68)   &5.56	    &2.677	    &0.02                   &0.37 \\ 
                        &$2^{+}$    &1.00                   &5.66	    &3.063	    &0.45                   &0.37 \\ 

        \multirow{2}{*}{${(nn\bar{s}\bar{b})}^{I=1}$}
                        &\multirow{2}{*}{$0^{+}$}
                                    &(0.53, 0.85)           &5.41	    &6.580	    &-                      &1.04, 0.69, 0.35 \\ 
                        &           &(-0.85, 0.53)          &5.41	    &6.204	    &-                      &1.04, 0.69, 0.35 \\ 
                        &\multirow{3}{*}{$1^{+}$}
                                    &(0.66, 0.75, 0.06)     &5.41	    &6.419	    &1.30, 0.49, -0.31      &1.04, 0.69, 0.35 \\ 
                        &           &(0.36, -0.38, 0.85)    &5.41	    &6.554	    &0.47, 0.26, 0.04       &1.04, 0.69, 0.35 \\ 
                        &           &(0.66, -0.55, -0.52)   &5.41	    &6.230	    &0.74, 0.20, -0.33      &1.04, 0.69, 0.35 \\ 
                        &$2^{+}$    &1.00                   &5.47	    &6.446	    &1.69, 0.64, -0.40      &1.05, 0.70, 0.36 \\ 
        \multirow{2}{*}{${(nn\bar{s}\bar{b})}^{I=0}$}
                        &\multirow{2}{*}{$0^{+}$}
                                    &(-0.66, -0.75)         &5.41	    &6.327	    &-                      &0.69 \\ 
                        &           &(-0.75, 0.66)          &5.41	    &5.988	    &-                      &0.69 \\ 
                        &\multirow{3}{*}{$1^{+}$}                            
                                    &(0.66, -0.23, 0.72)    &5.41	    &6.326	    &0.26                   &0.69 \\ 
                        &           &(0.46, 0.88, -0.14)    &5.41	    &6.454	    &0.47                   &0.69 \\ 
                        &           &(0.60, -0.42, -0.68)   &5.41	    &6.042	    &0.23                   &0.69 \\ 
                        &$2^{+}$    &1.00                   &5.48	    &6.431	    &0.64                   &0.70 \\ 
        \hline\hline
    \end{tabular}
\end{table*}

\renewcommand{\tabcolsep}{0.59cm}
\renewcommand{\arraystretch}{1.2}
\begin{table*}[!htbp]
    \caption{Predicted spectra of doubly heavy tetraquarks $nn\bar{c}\bar{c}$, $nn\bar{b}\bar{b}$ and $nn\bar{c}\bar{b}$. 
    Magnetic moments and charge radii are organized in the order of $I_{3}=1,\,0,\,-1$ for $I=1$.}
    \label{tab:2heavytetraquark1}
    \begin{tabular}{ccccccc}
        \hline\hline
        {\rm State}     &$J^{P}$    &{\rm Eigenvector}      &$R_{0}$    &$M_{bag}$  &$\mu_{bag}/\mu_{p}$    &$r_{E}$ \\ \hline
        \multirow{2}{*}{${(nn\bar{c}\bar{c})}^{I=1}$}
                        &\multirow{2}{*}{$0^{+}$}
                                    &(0.41, 0.91)           &5.24	    &4.343	    &-                      &0.54, 0.52, 0.92 \\ 
                        &           &(-0.91, 0.41)          &5.24	    &4.035	    &-                      &0.54, 0.52, 0.92 \\ 
                        &$1^{+}$    &1.00                   &5.22	    &4.117	    &0.49, -0.02, -0.51     &0.54, 0.52, 0.92 \\ 
                        &$2^{+}$    &1.00                   &5.32	    &4.179	    &1.00, -0.02, -1.04     &0.55, 0.53, 0.93 \\ 
        \multirow{2}{*}{${(nn\bar{c}\bar{c})}^{I=0}$}
                        &\multirow{2}{*}{$1^{+}$}
                                    &(-0.97, -0.24)         &5.24	    &3.926	    &-0.32                  &0.52 \\ 
                        &           &(-0.24, 0.97)          &5.24	    &4.205	    &0.29                   &0.52 \\ 
        \multirow{2}{*}{${(nn\bar{b}\bar{b})}^{I=1}$}
                        &\multirow{2}{*}{$0^{+}$}
                                    &(0.17, 0.98)           &4.83	    &11.093	    &-                      &0.91, 0.58, 0.38 \\ 
                        &           &(-0.98, 0.17)          &4.83	    &10.834	    &-                      &0.91, 0.58, 0.38 \\ 
                        &$1^{+}$    &1.00                   &4.83	    &10.854	    &0.65, 0.18, -0.28      &0.54, 0.52, 0.92 \\ 
                        &$2^{+}$    &1.00                   &4.88	    &10.878	    &1.31, 0.37, -0.56      &0.92, 0.59, 0.38 \\ 
        \multirow{2}{*}{${(nn\bar{b}\bar{b})}^{I=0}$}
                        &\multirow{2}{*}{$1^{+}$}
                                    &(-1.00, -0.08)         &4.83	    &10.655	    &0.06                   &0.58 \\ 
                        &           &(-0.08, 1.00)          &4.83	    &10.982	    &0.31                   &0.58 \\ 

        \multirow{2}{*}{${(nn\bar{c}\bar{b})}^{I=1}$}
                        &\multirow{2}{*}{$0^{+}$}
                                    &(0.31, 0.95)           &5.05	    &7.714	    &-                      &0.77, 0.24, 0.69 \\ 
                        &           &(-0.95, 0.31)          &5.05	    &7.440	    &-                      &0.77, 0.24, 0.69 \\ 
                        &\multirow{3}{*}{$1^{+}$}
                                    &(0.65, 0.76, 0.09)     &5.05	    &7.509	    &0.84, 0.08, -0.68      &0.77, 0.24, 0.69 \\ 
                        &           &(0.13, -0.22, 0.97)    &5.05	    &7.699	    &-0.06, -0.11, -0.17    &0.77, 0.24, 0.69 \\ 
                        &           &(0.75, -0.61, -0.24)   &5.05	    &7.466	    &0.94, 0.30, -0.34      &0.77, 0.24, 0.69 \\ 
                        &$2^{+}$    &1.00                   &5.12	    &7.531	    &1.16, 0.18, -0.80      &0.78, 0.25, 0.69 \\ 
        \multirow{2}{*}{${(nn\bar{c}\bar{b})}^{I=0}$}
                        &\multirow{2}{*}{$0^{+}$}
                                    &(-0.93, -0.36)         &5.05	    &7.502	    &-                      &0.24 \\ 
                        &           &(-0.36, 0.93)          &5.05	    &7.263	    &-                      &0.24 \\ 
                        &\multirow{3}{*}{$1^{+}$}
                                    &(0.91, -0.36, 0.21)    &5.05	    &7.519	    &0.20                   &0.24 \\ 
                        &           &(0.39, 0.92, -0.10)    &5.05	    &7.605	    &0.18                   &0.24 \\ 
                        &           &(-0.16, 0.17, 0.97)    &5.05	    &7.289	    &-0.12                  &0.24 \\ 
                        &$2^{+}$    &1.00                   &5.14	    &7.483	    &0.18                   &0.25 \\ 
        \hline\hline
    \end{tabular}
\end{table*}

\renewcommand{\tabcolsep}{0.83cm}
\renewcommand{\arraystretch}{1.2}
\begin{table*}[!htbp]
    \caption{Predicted spectra of doubly heavy tetraquarks $ss\bar{c}\bar{c}$, $ss\bar{b}\bar{b}$ and $ss\bar{c}\bar{b}$.}
    \label{tab:2heavytetraquark0}
    \begin{tabular}{ccccccc}
        \hline\hline
        {\rm State}     &$J^{P}$    &{\rm Eigenvector}      &$R_{0}$    &$M_{bag}$  &$\mu_{bag}/\mu_{p}$    &$r_{E}$ \\ \hline
        \multirow{2}{*}{$ss\bar{c}\bar{c}$}
                        &\multirow{2}{*}{$0^{+}$}
                                    &(-0.49, -0.87)         &5.30	    &4.520	    &-                      &0.89 \\ 
                        &           &(-0.87, 0.49)          &5.30	    &4.305	    &-                      &0.89 \\ 
                        &$1^{+}$    &1.00                   &5.30	    &4.382	    &-0.44                  &0.89 \\ 
                        &$2^{+}$    &1.00                   &5.39	    &4.433	    &-0.88                  &0.91 \\ 
        \multirow{2}{*}{$ss\bar{b}\bar{b}$}
                        &\multirow{2}{*}{$0^{+}$}
                                    &(-0.25, -0.97)         &4.92	    &11.231	    &-                      &0.31 \\ 
                        &           &(-0.97, 0.25)          &4.92	    &11.079	    &-                      &0.31 \\ 
                        &$1^{+}$    &1.00                   &4.94	    &11.098	    &-0.21                  &0.31 \\ 
                        &$2^{+}$    &1.00                   &4.98	    &11.119	    &-0.43                  &0.32 \\ 
        \multirow{2}{*}{$ss\bar{c}\bar{b}$}
                        &\multirow{2}{*}{$0^{+}$}
                                    &(-0.40, -0.92)         &5.13	    &7.873	    &-                      &0.65 \\ 
                        &           &(-0.92, 0.40)          &5.13	    &7.696	    &-                      &0.65 \\ 
                        &\multirow{3}{*}{$1^{+}$}
                                    &(0.70, 0.71, 0.11)     &5.13	    &7.759	    &-0.56                  &0.65 \\ 
                        &           &(0.17, -0.31, 0.93)    &5.13	    &7.858	    &-0.17                  &0.65 \\ 
                        &           &(0.70, -0.63, -0.34)   &5.13	    &7.720	    &-0.25                  &0.65 \\ 
                        &$2^{+}$    &1.00                   &5.20	    &7.779	    &-0.66                  &0.66 \\ 
        \hline\hline
    \end{tabular}
\end{table*}

\renewcommand{\tabcolsep}{0.73cm}
\renewcommand{\arraystretch}{1.3}
\begin{table*}[!htbp]
    \caption{Predicted spectra of doubly heavy tetraquarks $ns\bar{c}\bar{c}$, $ns\bar{b}\bar{b}$. 
    Magnetic moments and charge radii are organized in the order of $I_{3}=\frac{1}{2},\,-\frac{1}{2}$ for $I=\frac{1}{2}$.}
    \label{tab:2heavytetraquark1/2}
    \begin{tabular}{ccccccc}
        \hline\hline
        {\rm State}     &$J^{P}$    &{\rm Eigenvector}      &$R_{0}$    &$M_{bag}$  &$\mu_{bag}/\mu_{p}$    &$r_{E}$ \\ \hline
        \multirow{2}{*}{$ns\bar{c}\bar{c}$}
                        &\multirow{2}{*}{$0^{+}$}
                                    &(-0.44,-0.90)           &5.27	    &4.429	    &-                      &0.49, 0.91 \\ 
                        &           &(-0.89, 0.44)           &5.27	    &4.169	    &-                      &0.49, 0.91 \\ 
                        &\multirow{3}{*}{$1^{+}$}
                                    &(0.99, -0.07, 0.09)    &5.27	    &4.247	    &0.12, -0.48            &0.49, 0.91 \\ 
                        &           &(-0.11, -0.28, 0.95)   &5.27	    &4.314	    &0.31, -0.57            &0.49, 0.91 \\ 
                        &           &(-0.04, -0.96, -0.28)  &5.27	    &4.095	    &-0.35, -0.37           &0.49, 0.91 \\ 
                        &$2^{+}$    &1.00                   &5.36	    &4.305	    &0.07, -0.96            &0.50, 0.92 \\ 
        \multirow{2}{*}{$ns\bar{b}\bar{b}$}             
                        &\multirow{2}{*}{$0^{+}$}             
                                    &(-0.20, -0.98)         &4.88	    &11.160	    &-                      &0.61, 0.35 \\ 
                        &           &(-0.98, 0.20)          &4.88	    &10.956	    &-                      &0.61, 0.35 \\ 
                        &\multirow{3}{*}{$1^{+}$}
                                    &(1.00, -0.01, 0.03)    &4.88	    &10.975	    &0.24, -0.25            &0.61, 0.35 \\ 
                        &           &(-0.03, -0.09, 1.00)   &4.88	    &11.069	    &0.37, -0.55            &0.61, 0.35 \\ 
                        &           &(-0.01, -1.00, -0.09)  &4.88	    &10.814	    &0.05, 0.06             &0.61, 0.35 \\ 
                        &$2^{+}$    &1.00                   &4.93	    &10.997	    &0.45, -0.50            &0.62, 0.35 \\ 
        \hline\hline
    \end{tabular}
\end{table*}

\end{document}